\chapter{Дүгнэлт ба цаашдын зорилт}

\section{Дүгнэлт}
Энэхүү судалгаа болон төслийн хүрээнд дэлгүүрийн лангуун дээрх бараа бүтээгдэхүүнийг зурганд суурилан алдаагүй таньж, барааны тооллого болон аудитын шийдвэр гаргалтын системийг бүрэн автоматжуулах цогц ухаалаг системийг амжилттай хөгжүүлж туршлаа. Энэхүү систем нь хиймэл оюун ухааны объект таних алгоритмын тэргүүлэгч болох YOLOv8 загварыг, өндөр бүтээмжтэй FastAPI сервер болон хэрэглэгчдэд ээлтэй Flutter гар утасны аппликейшнтэй уялдуулж өгснөөрөө технологийн хувьд нэгдмэл, хурдан, оновчтой, хүртээмжтэй шийдэл болж чадсан юм. Туршилтын үр дүнгээр мэдрэлийн сүлжээний илрүүлэлтийн нарийвчлал болон серверийн хариу өгөх хурд өндөр байсан нь уг системийг бодит амьдрал дээр шууд хэрэгжүүлэх боломжтойг нотолж байна. Энэ нь жижиглэн худалдааны бизнес эрхлэгчид, түгээлтийн байгууллагуудад цаг хугацаа, хүний хөдөлмөрөө ихээхэн хэмнэх боломжийг бүрдүүлж байгаагаараа бодит ач холбогдолтой юм.

\section{Сул тал ба хязгаарлалт}
Туршилтын үед системийн хувьд зарим нэг хязгаарлалтууд ажиглагдсан. Үүнд, ижил төстэй хаяг шошготой боловч өөр төрлийн (жишээ нь, нэг ундааны үхэрийн нүдтэй болон анартай хувилбар зэрэг өнгө бага зэрэг ялгаатай) бараануудыг хольж алдаж таних тохиолдол зарим үед гарсан. Мөн зураг дарах үеийн гарны чичиргээ зэргээс шалтгаалан дүрс бүрсгэр (blur) гарсан тохиолдолд танилтын дундаж нарийвчлал унах магадлалтай. Түүнчлэн, загварыг зөвхөн хязгаарлагдмал тооны ангиллын бараан дээр сургасан нь одоогоор хязгаарлагдмал хүрээг дэмжихэд хүргэж байна.

\section{Цаашдын зорилт}
Цаашид энэхүү системийг дараах байдлаар өргөжүүлэн сайжруулах боломжтой:
\begin{itemize}
    \item Өгөгдлийн санг баяжуулж, зуу болон түүнээс дээш тооны олон ангиллын бараануудын хувьд детектор загварыг сургаж цар хүрээг нэмэгдүүлэх.
    \item Системийг зөвхөн гэрэл зургаар (Image) хязгаарлахгүйгээр шууд камерын бичлэгээс (Video Stream) эсвэл гар утсаа шилжүүлэх явцад (continuous scanning) шууд таньж анализ хийх боломжийг судлах.
    \item Олон утас зэрэг холбогдож ажиллахад зориулан серверийн ачааллыг оновчлох, мэдээллийн аюулгүй байдлыг ханган үүлэн технологид (Cloud infrastructure) байршуулж өргөтгөх боломжоор хангах.
\end{itemize}
